\begin{definition} \label{definition:symmetric_skew_symmetric_alternating_bilinear_form_on_a_module_over_a_ring}
    \TODO{AI generated, verify}
    Let $R$ be a \CrefAndHyperrefIfExist{definition:ring}{(not necessarily commutative) ring}, and let $\beta: M \times M \to R$ be a \CrefAndHyperrefIfExist{definition:multilinear_map_of_modules_over_rings}{bilinear form} on a two-sided \CrefAndHyperrefIfExist{definition:bimodule_over_rings_module_over_a_ring}{$R$-module} $M$.

    \begin{enumerate}
        \item $\beta$ is called \hldef{symmetric} if $\beta(m, n) = \beta(n, m)$ for all $m, n \in M$.

        \item $\beta$ is called \hldef{skew-symmetric} (or \hldef{antisymmetric}) if $\beta(m, n) = -\beta(n, m)$ for all $m, n \in M$.

        \item $\beta$ is called \hldef{alternating} if $\beta(m, m) = 0$ for all $m \in M$.
    \end{enumerate}

    Every alternating bilinear form is skew-symmetric. The converse holds whenever $2$ is invertible in $R$: if $\beta$ is skew-symmetric and $2$ is a unit, then $\beta(m,m) = -\beta(m,m)$ implies $\beta(m,m) = 0$.
\end{definition}
